% GENERATED FILE. Edit canonical lessons, then run npm run generate:books.
% canonical-source-hash: fnv1a64:7b8937c8c820afcf
% canonical-lessons: GE-C408-vorfahrt, GE-C408-ueberholen, GE-C408-draengeln, GE-C408-gegenverkehr, GE-C408-ausweichen

\chapter{Wer fährt zuerst}
\label{ch:ge-wer-fahrt-zuerst}
\hlchaptermodality{408}

\begin{chapteropening}
\textbf{By the end of this chapter:} \emph{I can talk about right of way and other drivers: right of way, overtaking, tailgating, oncoming traffic and swerving.}

Complete the last lesson of chapter 408: I can talk about right of way and other drivers: right of way, overtaking, tailgating, oncoming traffic and swerving.
\end{chapteropening}

\section[die Vorfahrt]{die Vorfahrt — the right of way}

\label{lesson:GE-C408-vorfahrt}



\begin{quote}
Before the new one: say the German for the fine, then the German for the penalty ticket.
\end{quote}

\subsection*{You'll want to know: die Vorfahrt}
\textbf{die Vorfahrt} — \textquotedblleft{}the right of way\textquotedblright{}.

\begin{quote}
An dieser Kreuzung hat der Bus Vorfahrt.
\end{quote}

\begin{quote}
\emph{At this junction the bus has right of way.}
\end{quote}

\emph{Vorfahrt achten} — give way, on a sign. In Austria: \emph{der Vorrang}; in Switzerland: \emph{der Vortritt}.

\begin{cousinweb}
\emph{Vorfahrt} is built from \emph{vor} 'before, ahead' and \emph{Fahrt} 'drive', from the inherited verb \emph{fahren} (EtymDB).
\end{cousinweb}

\begin{grammarlens}[title={What you've built}]
One more word for telling, explaining and describing at B1.
\end{grammarlens}

\subsection*{Your turn}
\begin{itemize}
\raggedright
  \item \emph{Say it:} \emph{die Vorfahrt}
  \item \emph{Say it:} \emph{die Vorfahrt}, once more
  \item \emph{Say it:} a sentence of your own with \emph{die Vorfahrt}
  \item \emph{Recall:} say the German for the fine, then the German for the penalty ticket, then say \emph{die Vorfahrt} again
\end{itemize}

\begin{tcolorbox}[breakable,colback=teal!4,colframe=teal!35!black,title={Before you move on}]
What does die Vorfahrt mean? (\textquotedblleft{}The right of way\textquotedblright{}.) Say it once more.
\end{tcolorbox}
\section[überholen]{überholen — to overtake; to overhaul}

\label{lesson:GE-C408-ueberholen}



\begin{quote}
Before the new one: say the German for the penalty ticket, then the German for the right of way.
\end{quote}

\subsection*{You'll want to know: überholen}
\textbf{überholen} — \textquotedblleft{}to overtake; to overhaul\textquotedblright{}.

\begin{quote}
Der LKW überholt uns auf der Landstraße.
\end{quote}

\begin{quote}
\emph{The lorry overtakes us on the country road.}
\end{quote}

Inseparable: \emph{er überholt}, \emph{er hat überholt}. \emph{Überholen} also means to overhaul a machine.

\begin{cousinweb}
\emph{überholen} is built from \emph{über} 'over, past' and \emph{holen} 'to fetch'; EtymDB lists Middle High German \emph{überholn}, which is usually said to have first meant fetching someone across a river.
\end{cousinweb}

\begin{grammarlens}[title={What you've built}]
One more word for telling, explaining and describing at B1.
\end{grammarlens}

\subsection*{Your turn}
\begin{itemize}
\raggedright
  \item \emph{Say it:} \emph{überholen}
  \item \emph{Say it:} \emph{überholen}, once more
  \item \emph{Say it:} a sentence of your own with \emph{überholen}
  \item \emph{Recall:} say the German for the penalty ticket, then the German for the right of way, then say \emph{überholen} again
\end{itemize}

\begin{tcolorbox}[breakable,colback=teal!4,colframe=teal!35!black,title={Before you move on}]
What does überholen mean? (\textquotedblleft{}To overtake; to overhaul\textquotedblright{}.) Say it once more.
\end{tcolorbox}
\section[drängeln]{drängeln — to tailgate; to push in, to jostle}

\label{lesson:GE-C408-draengeln}



\begin{quote}
Before the new one: say the German for the right of way, then the German for to overtake.
\end{quote}

\subsection*{You'll want to know: drängeln}
\textbf{drängeln} — \textquotedblleft{}to tailgate; to push in, to jostle\textquotedblright{}.

\begin{quote}
Hinter uns drängelt ein schwarzes Auto.
\end{quote}

\begin{quote}
\emph{A black car is tailgating us.}
\end{quote}

\emph{Drängeln} is also pushing in a queue: \emph{sich vordrängeln}.

\begin{cousinweb}
\emph{drängeln} is built on \emph{drängen} 'to push, to press', inherited from Old High German \emph{drengen} and Proto-Germanic \emph{þrangwijaną} (EtymDB). English \emph{throng} is usually counted as a cousin.
\end{cousinweb}

\begin{grammarlens}[title={What you've built}]
One more word for telling, explaining and describing at B1.
\end{grammarlens}

\subsection*{Your turn}
\begin{itemize}
\raggedright
  \item \emph{Say it:} \emph{drängeln}
  \item \emph{Say it:} \emph{drängeln}, once more
  \item \emph{Say it:} a sentence of your own with \emph{drängeln}
  \item \emph{Recall:} say the German for the right of way, then the German for to overtake, then say \emph{drängeln} again
\end{itemize}

\begin{tcolorbox}[breakable,colback=teal!4,colframe=teal!35!black,title={Before you move on}]
What does drängeln mean? (\textquotedblleft{}To tailgate; to push in, to jostle\textquotedblright{}.) Say it once more.
\end{tcolorbox}
\section[der Gegenverkehr]{der Gegenverkehr — the oncoming traffic}

\label{lesson:GE-C408-gegenverkehr}



\begin{quote}
Before the new one: say the German for to overtake, then the German for to tailgate.
\end{quote}

\subsection*{You'll want to know: der Gegenverkehr}
\textbf{der Gegenverkehr} — \textquotedblleft{}the oncoming traffic\textquotedblright{}.

\begin{quote}
In der Kurve kommt plötzlich Gegenverkehr.
\end{quote}

\begin{quote}
\emph{On the bend there is suddenly oncoming traffic.}
\end{quote}

No plural.

\begin{cousinweb}
\emph{Gegenverkehr} is built from \emph{gegen} 'against, towards' and \emph{Verkehr} 'traffic', from \emph{verkehren} 'to go to and fro', built on \emph{kehren} 'to turn'. The parts are inherited Germanic words.
\end{cousinweb}

\begin{grammarlens}[title={What you've built}]
One more word for telling, explaining and describing at B1.
\end{grammarlens}

\subsection*{Your turn}
\begin{itemize}
\raggedright
  \item \emph{Say it:} \emph{der Gegenverkehr}
  \item \emph{Say it:} \emph{der Gegenverkehr}, once more
  \item \emph{Say it:} a sentence of your own with \emph{der Gegenverkehr}
  \item \emph{Recall:} say the German for to overtake, then the German for to tailgate, then say \emph{der Gegenverkehr} again
\end{itemize}

\begin{tcolorbox}[breakable,colback=teal!4,colframe=teal!35!black,title={Before you move on}]
What does der Gegenverkehr mean? (\textquotedblleft{}The oncoming traffic\textquotedblright{}.) Say it once more.
\end{tcolorbox}
\section[ausweichen]{ausweichen — to swerve, to get out of the way; to avoid}

\label{lesson:GE-C408-ausweichen}



\begin{quote}
Before the new one: say the German for to tailgate, then the German for the oncoming traffic.
\end{quote}

\subsection*{You'll want to know: ausweichen}
\textbf{ausweichen} — \textquotedblleft{}to swerve, to get out of the way; to avoid\textquotedblright{}.

\begin{quote}
Die Fahrerin weicht einer Katze aus.
\end{quote}

\begin{quote}
\emph{The driver swerves to avoid a cat.}
\end{quote}

Separable, with the dative; Perfekt with \emph{sein}: \emph{sie ist der Katze ausgewichen}.

\begin{cousinweb}
\emph{ausweichen} is built from \emph{aus} 'out' and \emph{weichen} 'to give way, to yield', inherited from Old High German \emph{wīchan} and Proto-Germanic \emph{wīkwaną} (EtymDB).
\end{cousinweb}

\begin{grammarlens}[title={What you've built}]
One more word for telling, explaining and describing at B1.
\end{grammarlens}

\subsection*{Your turn}
\begin{itemize}
\raggedright
  \item \emph{Say it:} \emph{ausweichen}
  \item \emph{Say it:} \emph{ausweichen}, once more
  \item \emph{Say it:} a sentence of your own with \emph{ausweichen}
  \item \emph{Recall:} say the German for to tailgate, then the German for the oncoming traffic, then say \emph{ausweichen} again
\end{itemize}

\begin{tcolorbox}[breakable,colback=teal!4,colframe=teal!35!black,title={Before you move on}]
What does ausweichen mean? (\textquotedblleft{}To swerve, to get out of the way; to avoid\textquotedblright{}.) Say it once more.
\end{tcolorbox}
